\subsection*{The same local ink can ask four different questions}
The local question encoder loses distinctions between requests that share every centered nine-glyph neighborhood (Methods). We constructed four English questions about the same headword that differ in which fact and script the reader already knows and which two items are requested. Reordering the fields between clauses preserves every local window and its count. On 284 held-out headwords in four layouts the trained local writer returned the same stream for all four requests in every group. A larger control with global attention but no positional information did the same (Fig.~\ref{fig:scoped-request-writer}).

Giving the attention heads relative reading distance changed this behaviour. With familiar clause cues the ordered writer gave 4,396 of 4,544 exact answers. It answered all four requests correctly in 1,060 of 1,136 headword--layout groups (93.31 percentage points above unordered attention, 95\% card-clustered interval 90.49--95.95). It answered all 16 questions about 264 of the 284 headwords correctly. Panel c shows how the same local ink yields distinct source-grounded sentences for the headword \textit{\v{e}r} (ear). The two global models each had 694,152 parameters, with 181,761 trainable. Their ordered and unordered fits took 82.9 and 78.8 seconds, respectively, and each allocated less than 0.44\,GiB at peak.

With unfamiliar wording for both clauses and their fields, the ordered writer completed zero groups and only 561 of 4,544 single requests. Of its responses, 1,813 gave both the unwanted fact and the unwanted stage. A frozen factorial exchange isolated the wording effects. Changing only the words surrounding the fields reduced complete groups from 1,060 to 784. Changing only the requested-information cue reduced them to 246. Changing only the known-information cue reduced them to 16. Thus the position signal enabled clause assignment under familiar language; transfer across the two cue systems failed. Among the 96 independently authored requests per language the ordered writer gave 60 English and 56 Chinese answers, compared with 53 and 59 under unordered attention. Natural composition still yielded zero exact responses among 300 references requiring a new content glyph.
